% =====================================================================
%  说明书摘要
%
%  审查指南第一部分第一章 4.5.1：摘要文字部分应当写明
%      ① 发明的名称  ② 所属的技术领域  ③ 所要解决的技术问题
%      ④ 解决该问题的技术方案的要点  ⑤ 主要用途
%  五项缺一不可（「未写明发明名称或者不能反映技术方案要点的，应当通知
%  申请人补正」），且文字部分含标点不得超过 300 个字，不得使用标题，
%  不得使用商业性宣传用语。摘要中出现的附图标记应当加括号。
%
%  改动后务必重新数字数：node patent/tex/build.cjs --check
%  摘要附图必须是说明书附图中已有的一幅，且在请求书中写明图号。
% =====================================================================
\documentclass[12pt,UTF8,fontset=none]{ctexart}
\usepackage[abstract]{cnipa}

\begin{document}

% 摘要是单段，不编 [0001] 段号，因此用普通段落而不是 \p{}。
\cnabstract{本发明公开一种大语言模型智能体的工具调用控制方法及系统，属计算机软件技术领域。要解决的技术问题是：智能体反复发起不产生新观察的工具调用而空耗算力与令牌，以调用是否重复为唯一判据既误杀轮询又漏判入参微变的同义调用。技术方案要点：对每次已完成调用生成调用指纹，由剔除时间戳等易变字段并规范化后的输出生成证据指纹；外置为产物文件的输出以内容摘要值而非路径为依据；维护证据指纹集合与无新证据连续计数，新证据清零，重复或空累加；达第一阈值注入策略变更指示，达第二阈值且两项计数分别达第三、第四阈值时执行前拒绝；上下文压缩时由已存储的调用—结果配对重建上述状态。可用于智能体运行时，减少无效工具调用与网络请求。}

\end{document}
